% This document is created to create questions and solutions for a Math Exam. The solutions would be on the left side of the solution portion, the Math step by step, and on the right side, it will be the explanation of what the Math is doing in words.

\documentclass{article}

\usepackage{graphicx} % Required for inserting images
\usepackage{amsmath, amssymb, amsthm} % For better use of Math expressions
\usepackage{hyperref} % for adding hyperlinks to the document. 
\usepackage{comment}
\usepackage{marginnote} % Adding notes to the margin of the document
\reversemarginpar % Optional: to switch margin sides
\usepackage{xcolor} % For creating the the colorbox
% For creating tcolorboxes. 
\usepackage{tikz}
\usepackage[most]{tcolorbox}
\usepackage{lipsum}
\usepackage{enumitem} % for enumerating items with letters from the alphabet.
\usepackage{tabularx} % Enables tables with flexible-with columns
\usepackage{etoolbox} % For custom labels and logic.
\usepackage{array} % Extend column definitions in tables
\usepackage{witharrows} % for adding arrows

% ==== Customization ====
\usepackage [letterpaper, top=1.0in, bottom=1.0in, left=1.0in, right=1.0in, heightrounded]{geometry} % Margins of the document. 
\renewcommand{\baselinestretch}{1.2} % space line length
\setlength{\parindent}{0pt} % Change the size of the tab
\setlength{\parskip}{0.8em} %Change the size between paragraphs

\title{\textbf{Cal 2 - Exam Questions - Solution with Steps and SubSteps }}
\author{}
\date{}


% ==== Custom Macros ====
% Create a general counter that resets automatically per question
\newcounter{solutionMainStep}
\newcounter{solutionSubStep}[solutionMainStep]

% Define a macro that takes question number, summary of step on (left), and explanation of the step on (right). 
\renewcommand{\thesolutionSubStep}{\arabic{solutionMainStep}.\arabic{solutionSubStep}}

% Macro for main step 
\newcommand{\mainStep}[2]{%
    \stepcounter{solutionMainStep}
    \setcounter{solutionSubStep}{0} % reset sub steps.
    \label{a-#1:\arabic{solutionMainStep}}% adding a label to the macro for later reference.
    % Code to organize and display the bubbles
    \begin{tcolorbox}[
        breakable,
        colback = gray!5!white, 
        colframe = gray!80!black, 
        boxrule = 0.5pt,
        arc = 2mm, 
        left=0pt, right=0pt, top=2pt, bottom=2pt,
        before skip=2pt, after skip = 2pt,
        enhanced]
        \textbf{Step \arabic{solutionMainStep}:} \textbf{\small #2}
    \end{tcolorbox}
}

% Macro for substep (e.g., "Step 1.1:")
\newcommand{\solutionStep}[3]{%
    \stepcounter{solutionSubStep}
    \label{a-#1:\arabic{solutionMainStep}.\arabic{solutionSubStep}}% adding a label to the macro for later reference.
    % Code to organize and display the bubbles
    \begin{tcolorbox}[
        breakable,
        colback = gray!5!white,
        colframe = gray!80!black,
        boxrule = 0.5pt,
        arc = 2mm,
        left=0pt, right=0pt, top=2pt, bottom=2pt,
        before skip=2pt, after skip=2pt,
        enhanced]
        \begin{tabularx}{\textwidth}{>{\raggedright\arraybackslash}p{5cm} X}
            \textbf{Step \arabic{solutionMainStep}.\arabic{solutionSubStep}:}\\
            \textbf{\small #2} & #3
        \end{tabularx}
    \end{tcolorbox}
}

% To reference a step:
% Use \ref{a-<question>:<step>} anywhere in the text.
% Example: \ref{a-2:1.2} refers to Step 1.2 of Question 2.

% ==== Beginning of Document ==== 
\begin{document}
\maketitle
% QUESTION: 1
% STATUS: REVIEWED
% AUTHOR: AJ
% CREATED: 
% LAST UPDATED: 
% NOTES: 

\section*{Question 1}
    Let $f(x) = x^5 + x^3 - 4$. Find the value of $\frac{d}{dx}f^{-1}(x)$ at the point $x = -2$.  
    
    % Possible answers
    \begin{enumerate}[label = \alph*.]
        \item $-\frac{1}{44}$
        \item $\frac{1}{92}$
        \item $\frac{1}{8}$
        \item I do not know
    \end{enumerate}

\setcounter{solutionMainStep}{0}
\setcounter{solutionSubStep}{0}

% Beginning of Solution Box
%\begin{tcolorbox}[colback= blue!5!white, colframe= blue!50!black, title=\textbf{Solution: c) $3 \ln |x^5 + 15| + C $}]
 \begin{tcolorbox}[breakable, colback= blue!5!white, colframe= blue!50!black, title=\textbf{Solution: c) $\frac{1}{8}$}]

    % Step-by-step modular solution.
    %This solution is meant to be long with a lot more steps so that a 5th grader can understand the solution.

    % ------- Start here adding the step-by-step solution --------

    % === Step Formatting Summary ===
    % \mainStep{<question#>}{<Step title>}
    %   Creates a main step like "Step 1:", resets substep counter.
    %
    % \solutionStep{<question#>}{<Short title>}{<Detailed explanation>}
    %   Creates a substep like "Step 1.1:" under the current main step.

    % Step 1
     \mainStep{1}{Find $f'(x)$(derivative of $f(x)$).}

    % Step 1.1
    \solutionStep{1}{Apply the power rule to $f'(x)$}{\colorbox{yellow}{Power Rule:} if $g(x) = x^n$, then $g'(x) = nx^{n-1}$. }

    % Step 1.2
    \solutionStep{1}{Differentiate $x^5$}{ $\frac{d}{dx}(x^5)$ = $5x^4$  }

    % Step 1.3
    \solutionStep{1}{ Differentiate $x^3$}{ $\frac{d}{dx}(x^3)$ = $3x^2$  }

     % Step 1.4
    \solutionStep{1}{Differentiate the constant -4}{ $\frac{d}{dx}(-4)$ = $0$  }

     % Step 1.5
    \solutionStep{1}{Combine the results}{ $f'(x) = 5x^4+3x^2$  }

    % Step 2
    \mainStep{1}{Find $f^{-1}(-2)$ by using \colorbox{yellow} {Inverse Function}: if $f(a)=b$ then $f^{-1}(b) =a$}

    % Step 2.1
    \solutionStep{1}{Set up the equation $f(x) = -2.$}{$x^5+x^3-4 = -2$}

    % Step 2.2 Version 1
    \solutionStep{1}{Solve for x: }{

    \begin{flushleft}
        \begin{tabbing}
        \hspace{4.5cm} \= \kill % Set tab stop for aligning explanations to the right.
            % Start math block and align all equations at the equal sign.
            $\begin{aligned}[t] % [t] use to align the top of this aligned environment with the top of the line or cell it's in
                x^5 + x^3 - 4 &= -2 \\
                x^5 + x^3 - 4 \colorbox{yellow}{+ 4}  &= -2 \colorbox{yellow}{+ 4} \\
                x^5 + x^3 &= 2
            \end{aligned}$ \>
            % Now list matching explanations, one per line
            \begin{tabular}[t]{@{}l@{}} % Left-align text with no extra padding
                Original equation. \\ % English explanations
                Add 4 to each side. \\
                Simplify both sides.
            \end{tabular}
        \end{tabbing}
    \end{flushleft}
    }

    % Step 2.3
    \solutionStep{1}{Test simple value 1 to find solution}{
    \begin{flushleft}
        \begin{tabbing}
        \hspace{4.5cm} \= \kill % Set tab stop for aligning explanations to the right.
            % Start math block and align all equations at the equal sign.
            $\begin{aligned}[t] % [t] use to align the top of this aligned environment with the top of the line or cell it's in
                f(x) &= x^5 + x^3 - 4 \\
                f(\colorbox{yellow}{1}) &= \colorbox{yellow}{1}^5 + \colorbox{yellow}{1}^3 - 4  \\
                f(1)&= 1 + 1 - 4 \\
                f(1) &= -2
            \end{aligned}$ \>
            % Now list matching explanations, one per line
            \begin{tabular}[t]{@{}l@{}} % Left-align text with no extra padding
                \\ [1ex]% English explanations
                Try $x=1$ in $f(x)$. \\ [1ex]
                Simplify. \\ [1ex]

            \end{tabular}
        \end{tabbing}
    \end{flushleft}
    }

    % Step 2.4
    \solutionStep{1}{Conclude}{Since $f(1) = -2$,  $f^{-1}(-2) = 1$ due to Inverse Function}

     % Step 3
    \mainStep{1}{ Evaluate $f'$ at $f^{-1}(-2)$}

    % Step 3.1
    \solutionStep{1}{Substitute $x=1$ into $f'(x)$}{
    \begin{flushleft}
        \begin{tabbing}
        \hspace{4.5cm} \= \kill % Set tab stop for aligning explanations to the right.
            % Start math block and align all equations at the equal sign.
            $\begin{aligned}[t] % [t] use to align the top of this aligned environment with the top of the line or cell it's in
                 f'(x) &= 5x^4+3x^2\\
                 f'(\colorbox{yellow}{1}) &= 5(\colorbox{yellow}{1})^4 + 3(\colorbox{yellow}{1})^2 \\
                 f'(1) &= 5(1) + 3(1)\\
                 f'(1) &= 5 + 3 \\
                 f'(1) &= 8
            \end{aligned}$ \>
            % Now list matching explanations, one per line
            \begin{tabular}[t]{@{}l@{}} % Left-align text with no extra padding
                \\ [1ex]% English explanations
                Substitute $x$ with $1$ in $f'(x)$ \\
                \\
                \\

            \end{tabular}
        \end{tabbing}
    \end{flushleft}
    }

     % Step 4
    \mainStep{1}{Apply the inverse function derivative formula}

    % Step 4.1
    \solutionStep{1}{Substitute into the formula}{
    \begin{flushleft}
        \begin{tabbing}
        \hspace{4.5cm} \= \kill % Set tab stop for aligning explanations to the right.
            % Start math block and align all equations at the equal sign.
            $\begin{aligned}[t] % [t] use to align the top of this aligned environment with the top of the line or cell it's in
                 \frac{d}{dx}f^{-1}(x) &= \frac{1}{f'(f^{-1}(x))}\\
                 \frac{d}{dx}f^{-1}(x) &= \frac{1}{f'(f^{-1}(\colorbox{yellow}{-2}))} \\
                 \frac{d}{dx}f^{-1}(x) &= \frac{1}{\colorbox{yellow}{f'(1)}}\\
                \frac{d}{dx}f^{-1}(x) &= \frac{1}{8}
            \end{aligned}$ \>
            % Now list matching explanations, one per line
            \begin{tabular}[t]{@{}l@{}} % Left-align text with no extra padding
                Original Problem Equation.\\ [2.5ex]% English explanations
                Substitute $x$ with $-2$\\ [2.5ex]
                 $f'(1) = 8$\\ [2.5ex]
            \end{tabular}
        \end{tabbing}
    \end{flushleft}
    }

     % Step 5
    \mainStep{1}{Final Answer: the correct option is C. $\frac{1}{8}$ }
    
\end{tcolorbox}
\newpage
% QUESTION: 2
% STATUS: DRAFT
% AUTHOR: DG
% CREATED: 2026-02-02
% LAST UPDATED: 2026-02-02
% NOTES: 

% Qustion
\section*{Question 2}
Evaluate $\int\frac{15x^4}{x^5+15}\,dx$.  
    
% Multiple Choices
\begin{enumerate}[label = \alph*.]
    \item $\frac{3x^5}{\frac{1}{6}x^6+15x}+C$
    \item $15\ln|x|+\frac{1}{5}x^4+C$
    \item $3\ln|x^5+15|+C$
    \item I do not know
\end{enumerate}

\setcounter{solutionMainStep}{0}
\setcounter{solutionSubStep}{0}

% Solution Box
\begin{tcolorbox}[breakable, colback= blue!5!white, colframe= blue!50!black, title=\textbf{Solution: c) $3\ln|x^5+15|+C$}]

    % Step 1
    \mainStep{2}{Use U Substitution, Let $ u = x^5+15$}

    % Step 1.1
    \solutionStep{2}{Differentiate both sides of the equation }{ $\frac{du}{dx} = \frac{d}{dx}(x^5+15)$  }

    % Step 1.2
    \solutionStep{2}{Differentiate $x^5$}{ $\frac{d}{dx}(x^5) = 5x^4$  }
        
    % Step 1.3
    \solutionStep{2}{Differentiate the constant $15$}{ $\frac{d}{dx}(15) = 0$  }

    % Step 1.4
    \solutionStep{2}{Combine the results}{ $ \frac{d}{dx}(x^5+15) = 5x^4$  }

    % Step 1.5
    \solutionStep{2}{Multiply both sides by $dx$}{ $ du = 5x^4 dx$  }

    % Step 2
    \mainStep{2}{Rewrite $\int\frac{15x^4}{x^5+15}\,dx$ in terms of u.}

    % Step 2.1
    \solutionStep{2}{Rewrite numerator}{$15x^4dx = 3(5x^4dx) = 3du$ since $ du = 5x^4 dx$}

    % Step 2.2 
    \solutionStep{2}{Rewrite denominator}{$x^5+15 = u$}

    % Step 2.3
    \solutionStep{2}{Combine the fraction}{$\int\frac{15x^4}{x^5+15}\,dx = \int\frac{3}{u}du$}

    % Step 3
    \mainStep{2}{ Evaluate $\int\frac{3}{u}du$}
        
    % Step 3.1
    \solutionStep{2}{Factor out the constant}{$\int\frac{3}{u}du = 3\int\frac{1}{u}du$}

    % Step 3.2
    \solutionStep{2}{Evaluate the integral}{$\int\frac{1}{u}du = \ln{|u|} + C$}
        
    % Step 3.3
    \solutionStep{2}{Combine the solution}{$\int\frac{3}{u}du = 3\int\frac{1}{u}du = 3(\ln{|u|})+ C = 3\ln{|u|} + C$}
        
    % Step 4
    \mainStep{2}{Substitute back to x}

    % Step 4.1
    \solutionStep{2}{Replace $u$ with $x^5+15$}{$3\ln{|u|} + C =3\ln{|x^5+15|} + C$}

    % Step 5
    \mainStep{2}{Final Answer: the correct option is C. $3\ln|x^5+15|+C$ }
\end{tcolorbox}
\newpage  
% QUESTION: 3
% STATUS: DRAFT
% AUTHOR: DG
% CREATED: 2026-02-02
% LAST UPDATED: 2026-02-03
% NOTES: 

% Qustion
\section*{Question 3}
Evaluate $\frac{d}{dx}(\sin x)^{\ln x}$.  
    
% Multiple Choices
\begin{enumerate}[label = \alph*.]
    \item $(\ln x)(\sin x)^{(\ln x)-1}\cdot(\cos x)$
    \item $(\sin x)^{\ln x}\cdot\ln(\sin x)\cdot\frac{1}{x}$
    \item $(\sin x)^{\ln x}\left(\frac{1}{x}\cdot\ln(\sin x)+\ln x\cdot\cot x\right)$
    \item I do not know
\end{enumerate}

\setcounter{solutionMainStep}{0}
\setcounter{solutionSubStep}{0}

% Solution Box
\begin{tcolorbox}[breakable, colback= blue!5!white, colframe= blue!50!black, title=\textbf{Solution: c) $(\sin x)^{\ln x}\left(\frac{1}{x}\cdot\ln(\sin x)+\ln x\cdot\cot x\right)$}]

    % Step 1
    \mainStep{3}{Use logarithmic differentiation}

    % Step 1.1
    \solutionStep{3}{Rewrite the expression as a variable}{$y = (\sin x)^{\ln x}$}

    % Step 1.2
    \solutionStep{3}{Take natural log on both sides}{$\ln (y) = \ln ((\sin x)^{\ln x})$}
        
    % Step 1.3
    \solutionStep{3}{Apply the log power rule}{\colorbox{yellow}{Log Power Rule:} $\ln (a^b) = b\ln (a)$. }

    % Step 1.4
    \solutionStep{3}{Find $\ln ((\sin x)^{\ln x})$}{ $\ln ((\sin x)^{\ln x}) = (\ln x) \ln {(\sin x)}$  }

    % Step 2
    \mainStep{3}{Differentiate both sides of the equation}

    % Step 2.1
    \solutionStep{3}{Differentiate $\ln y$ using chain rule}{\colorbox{yellow}{Chain Rule:} if functions $f$ and $g$ are both differentiable and $F(x) = f(g(x))$, then $F'(x) = f'(g(x))g(x)$}

    % Step 2.2 
    \solutionStep{3}{Apply the chain rule}{$\frac{d}{dx}(\ln y) = (\ln y)'(y)'$}

    % Step 2.3
    \solutionStep{3}{Evaluate $(\ln y)'(y)'$}{$(\ln y)'= \frac{1}{y}$ and $y' = \frac{dy}{dx}$, so $(\ln y)'(y)' = \frac{1}{y} \frac{dy}{dx}$}

    % Step 2.4
    \solutionStep{3}{Differentiate $(\ln x) \ln {(\sin x)}$ using product rule}{\colorbox{yellow}{Product Rule:} if functions $f$ and $g$ are both differentiable and $F(x) = f(x)g(x)$, then $F'(x) = f(x)g'(x) + f'(x)g(x)$}

     % Step 2.5 
    \solutionStep{3}{Apply the product rule}{$\frac{d}{dx}(\ln x) \ln {(\sin x)} = (\ln x)' \ln {(\sin x)} + (\ln x) \ln {(\sin x)}'$}

    % Step 3
    \mainStep{3}{Evaluate $(\ln x)' \ln {(\sin x)} + (\ln x) \ln {(\sin x)}'$}
        
    % Step 3.1
    \solutionStep{3}{Differentiate $\ln x$}{$(\ln x)'=\frac{1}{x}$}

    % Step 3.2
    \solutionStep{3}{Differentiate $\ln (\sin x)$ using chain rule}{$(\ln (\sin x))'=(\ln '(\sin x))(\sin x)'$}

    
    % Step 3.3
    \solutionStep{3}{Differentiate $\sin x$}{$(\sin x)' = \cos x$}
    
    % Step 3.4
    \solutionStep{3}{Combine the results}{$(\ln(\sin x))' = \frac{\cos x}{\sin x} = \cot x$}
    
    % Step 3.5
    \solutionStep{3}{Substitute derivatives back into the product rule}{$\frac{d}{dx}\big[(\ln x)\ln(\sin x)\big] = \frac{1}{x}\ln(\sin x)+(\ln x)\cot x$}
    
    % Step 4
    \mainStep{3}{Solve for $\dfrac{dy}{dx}$}
    
    % Step 4.1
    \solutionStep{3}{Set both derivatives equal}{$\frac{1}{y}\frac{dy}{dx} = \frac{1}{x}\ln(\sin x)+(\ln x)\cot x$}
    
    % Step 4.2
    \solutionStep{3}{Multiply both sides by $y$}{$\frac{dy}{dx} = y\left(\frac{1}{x}\ln(\sin x)+(\ln x)\cot x\right)$}
    
    % Step 4.3
    \solutionStep{3}{Replace $y$ with $(\sin x)^{\ln x}$}{$\frac{d}{dx}(\sin x)^{\ln x} = (\sin x)^{\ln x}\left(\frac{1}{x}\ln(\sin x)+(\ln x)\cot x\right)$}
    
    % Step 5
    \mainStep{3}{Final Answer: the correct option is C. $(\sin x)^{\ln x}\left(\frac{1}{x}\cdot\ln(\sin x)+\ln x\cdot\cot x\right)$ }

   

\end{tcolorbox}
\newpage  



% QUESTION: 4
% STATUS: DRAFT
% AUTHOR: DG
% CREATED: 2026-02-02
% LAST UPDATED: 2026-02-07
% NOTES: 
% hightlight the operatiosn we are doing in the equation

% Question
\section*{Question 4}
Evaluate $\displaystyle\int x^3\cdot 5^{x^4}\,dx$.
    
% Multiple Choices
\begin{enumerate}[label = \alph*.]
    \item $\displaystyle\frac{5^{x^4}\ln 5}{4}+C$
    \item $\displaystyle\frac{x^4\cdot 5^{x^4}}{4}+C$
    \item $\displaystyle\frac{5^{x^4}}{4\ln 5}+C$
    \item I do not know
\end{enumerate}

\setcounter{solutionMainStep}{0}
\setcounter{solutionSubStep}{0}

% Solution Box
\begin{tcolorbox}[breakable, colback= blue!5!white, colframe= blue!50!black, title=\textbf{Solution: c) $\displaystyle\frac{5^{x^4}}{4\ln 5}+C$}]

    % Step 1
    \mainStep{4}{Use U-Substitution, Let $u = x^4$}

    % Step 1.1
    \solutionStep{4}{Differentiate both sides of the equation}{ $\frac{du}{dx} = \frac{d}{dx}(x^4)$ }

    % Step 1.2
    \solutionStep{4}{Apply the power rule}{\colorbox{yellow}{Power Rule:} if $g(x) = x^n$, then $g'(x) = nx^{n-1}$. }

    % Step 1.3
    \solutionStep{4}{Differentiate $x^4$}{ $\frac{d}{dx}(x^4) = 4x^{4-1} = 4x^3$ }

    % Step 1.4
    \solutionStep{4}{Combine the results}{ $\frac{du}{dx} = 4x^3$ }

    % Step 2
    \mainStep{4}{Convert the equation into differential form}

    % Step 2.1
    \solutionStep{4}{Multiply both sides by $dx$}{ $du = 4x^3 dx$ }

    % Step 2.2
    \solutionStep{4}{Solve for $x^3 dx$}{ Divide both sides by 4: $\displaystyle \frac{1}{4}du = x^3 dx$ }

    % Step 3
    \mainStep{4}{Rewrite the integral in terms of $u$}

    % Step 3.1
    \solutionStep{4}{Substitute $u = x^4$}{ $5^{x^4} = 5^u$ }

    % Step 3.2
    \solutionStep{4}{Substitute $x^3 dx = \frac{1}{4}du$}{ $\displaystyle\int x^3\cdot 5^{x^4}\,dx = \int 5^u \cdot \frac{1}{4}du$ }

    % Step 4
    \mainStep{4}{Evaluate $\int 5^u \cdot \frac{1}{4}du$}

    % Step 4.1
    \solutionStep{4}{Use constant multiple rule}{\colorbox{yellow}{Constant Multiple Rule:} $\int c f(x)dx = c\int f(x)dx$. }

    % Step 4.2
    \solutionStep{4}{Apply the rule with $c = \frac{1}{4}$}{ $\int 5^u \cdot \frac{1}{4}du = \frac{1}{4}\int 5^u du$ }

    % Step 4.3
    \solutionStep{4}{Use exponential integral rule}{\colorbox{yellow}{Exponential Integral Rule:} $\int a^u du = \frac{a^u}{\ln a} + C$. }

    % Step 4.4
    \solutionStep{4}{Apply the rule with $a=5$}{ $\int 5^u du = \frac{5^u}{\ln 5} + C$ }

    % Step 4.5
    \solutionStep{4}{Multiply by $\frac{1}{4}$}{ $\frac{1}{4}\left(\frac{5^u}{\ln 5} + C\right) = \frac{5^u}{4\ln 5} + C$ }

    % Step 5
    \mainStep{4}{Substitute back to $x$}

    % Step 5.1
    \solutionStep{4}{Replace $u$ with $x^4$}{ $\frac{5^u}{4\ln 5} + C = \frac{5^{x^4}}{4\ln 5} + C$ }

    % Step 6
    \mainStep{4}{Final Answer: the correct option is C. $\displaystyle\frac{5^{x^4}}{4\ln 5}+C$ }

\end{tcolorbox}
\newpage

% QUESTION: 5
% STATUS: DRAFT
% AUTHOR: DG
% CREATED: 2026-02-02
% LAST UPDATED: 2026-02-07
% NOTES: 
% hightlight the operatiosn we are doing in the equation

% Question
\section*{Question 5}
Evaluate $\displaystyle\int\frac{(\log_8 x)^2}{x}\,dx$.
    
% Multiple Choices
\begin{enumerate}[label = \alph*.]
    \item $\displaystyle(\ln 8)\cdot(\log_8 x)^2+C$
    \item $\displaystyle\frac{\ln 8}{3}\cdot(\log_8 x)^3+C$
    \item $\displaystyle\frac{(\ln 8)^2}{3}\cdot(\log_8 x)^3+C$
    \item I do not know
\end{enumerate}

\setcounter{solutionMainStep}{0}
\setcounter{solutionSubStep}{0}

% Solution Box
\begin{tcolorbox}[breakable, colback= blue!5!white, colframe= blue!50!black, title=\textbf{Solution: b) $\displaystyle\frac{\ln 8}{3}\cdot(\log_8 x)^3+C$}]

    % Step 1
    \mainStep{5}{Use U-Substitution, Let $u = \log_8 x$}

    % Step 1.1
    \solutionStep{5}{Rewrite $\log_8 x$ using change of base formula}{\colorbox{yellow}{Change of Base Formula:} $\log_b x = \frac{\ln x}{\ln b}$. }

    % Step 1.2
    \solutionStep{5}{Apply the formula with $b=8$}{ $u = \log_8 x = \frac{\ln x}{\ln 8}$ }

    % Step 2
    \mainStep{5}{Find $du$}

    % Step 2.1
    \solutionStep{5}{Differentiate both sides of the equation}{ $\frac{du}{dx} = \frac{d}{dx}\left(\frac{\ln x}{\ln 8}\right)$ }

    % Step 2.2
    \solutionStep{5}{Use constant multiple rule for derivatives}{\colorbox{yellow}{Constant Multiple Rule:} $\frac{d}{dx}[c f(x)] = c f'(x)$. }

    % Step 2.2
    \solutionStep{5}{Apply the formula with c = $c = ln 8$}{ $\frac{d}{dx}\frac{\ln x}{ln 8} = \frac{1}{ln 8}\frac{d}{dx}(ln x)$ }

    % Step 2.3
    \solutionStep{5}{Differentiate $\ln x$}{ $\frac{d}{dx}(\ln x) = \frac{1}{x}$ }

    % Step 2.4
    \solutionStep{5}{Combine the results}{ $\frac{du}{dx} = \frac{1}{\ln 8}\cdot\frac{1}{x} = \frac{1}{x\ln 8}$ }

    % Step 3
    \mainStep{5}{Convert into differential form}

    % Step 3.1
    \solutionStep{5}{Multiply both sides by $dx$}{ $du = \frac{1}{x\ln 8}dx$ }

    % Step 3.2
    \solutionStep{5}{Solve for $\frac{1}{x}dx$}{ Multiply both sides by $\ln 8$: $(\ln 8)du = \frac{1}{x}dx$ }

    % Step 4
    \mainStep{5}{Rewrite the integral in terms of $u$}

    % Step 4.1
    \solutionStep{5}{Substitute $u = \log_8 x$}{ $(\log_8 x)^2 = u^2$ }

    % Step 4.2
    \solutionStep{5}{Substitute $\frac{1}{x}dx = (\ln 8)du$}{ $\displaystyle\int\frac{(\log_8 x)^2}{x}dx = \int u^2 (\ln 8)du$ }
    
    % Step 4.3
    \solutionStep{5}{Use constant multiple rule}{\colorbox{yellow}{Constant Multiple Rule:} $\int c f(x)dx = c\int f(x)dx$. }

    % Step 4.4
    \solutionStep{5}{Rewrite the integral}{ $\int u^2 (\ln 8)du = (\ln 8)\int u^2 du$ }

    % Step 5
    \mainStep{5}{Evaluate $(\ln 8)\int u^2 du$}

    % Step 5.1
    \solutionStep{5}{Apply power rule for integrals}{\colorbox{yellow}{Power Rule for Integrals:} if $n\neq -1$, $\int u^n du = \frac{u^{n+1}}{n+1}+C$. }

    % Step 5.2
    \solutionStep{5}{Apply the rule with $n=2$}{ $\int u^2 du = \frac{u^3}{3}+C$ }

    % Step 5.3
    \solutionStep{5}{Multiply by $\ln 8$}{ $(\ln 8)\left(\frac{u^3}{3}+C\right) = \frac{\ln 8}{3}u^3 + C$ }

    % Step 6
    \mainStep{5}{Substitute back to $x$}

    % Step 6.1
    \solutionStep{5}{Replace $u$ with $\log_8 x$}{ $\frac{\ln 8}{3}u^3 + C = \frac{\ln 8}{3}(\log_8 x)^3 + C$ }

    % Step 7
    \mainStep{5}{Final Answer: the correct option is B. $\displaystyle\frac{\ln 8}{3}\cdot(\log_8 x)^3+C$ }

\end{tcolorbox}
\newpage

% QUESTION: 6
% STATUS: DRAFT
% AUTHOR: DG
% CREATED: 2026-02-02
% LAST UPDATED: 2026-02-07
% NOTES: Need to improve the following
%1. Need to expand Step 1.1 and 1.2
%2. More steps for 3.6 and 4.5
%3. Show the details of how to to go from 5.1 to 5.2

% Question
\section*{Question 6}
Evaluate $\displaystyle\lim_{x\to\infty}\frac{\ln(3x^2+7)}{\ln(5x^2+6)}$.
    
% Multiple Choices
\begin{enumerate}[label = \alph*.]
    \item $\displaystyle\ln\left(\frac{3}{5}\right)$
    \item 1
    \item $\displaystyle\frac{5}{3}$
    \item I do not know
\end{enumerate}

\setcounter{solutionMainStep}{0}
\setcounter{solutionSubStep}{0}

% Solution Box
\begin{tcolorbox}[breakable, colback= blue!5!white, colframe= blue!50!black, title=\textbf{Solution: b) 1}]

    % Step 1
    \mainStep{6}{Identify the limit form}

    % Step 1.1
    \solutionStep{6}{Evaluate the numerator as $x\to\infty$}{ $3x^2+7 \to \infty$, therefore $\ln(3x^2+7)\to\infty$ }

    % Step 1.2
    \solutionStep{6}{Evaluate the denominator as $x\to\infty$}{ $5x^2+6 \to \infty$, therefore $\ln(5x^2+6)\to\infty$ }

    % Step 1.3
    \solutionStep{6}{Determine the indeterminate form}{ The limit is of the form $\frac{\infty}{\infty}$ }

    % Step 2
    \mainStep{6}{Apply L'Hôpital's Rule}

    % Step 2.1
    \solutionStep{6}{State L'Hôpital's Rule}{\colorbox{yellow}{L'Hôpital's Rule:} If $\lim_{x\to a}\frac{f(x)}{g(x)}$ is $\frac{0}{0}$ or $\frac{\infty}{\infty}$, then $\lim_{x\to a}\frac{f(x)}{g(x)} = \lim_{x\to a}\frac{f'(x)}{g'(x)}$ }

    % Step 2.2
    \solutionStep{6}{Differentiate numerator and denominator}{ $\displaystyle\lim_{x\to\infty}\frac{\ln(3x^2+7)}{\ln(5x^2+6)} = \lim_{x\to\infty}\frac{\frac{d}{dx}\ln(3x^2+7)}{\frac{d}{dx}\ln(5x^2+6)}$ }

    % Step 3
    \mainStep{6}{Differentiate the numerator}

    % Step 3.1
    \solutionStep{6}{Use chain rule}{\colorbox{yellow}{Chain Rule:} if $F(x)=\ln(h(x))$, then $F'(x)=\frac{h'(x)}{h(x)}$. }

    % Step 3.2
    \solutionStep{6}{Differentiate inside function $3x^2+7$}{\colorbox{yellow}{Power Rule:} if $g(x)=x^n$, then $g'(x)=nx^{n-1}$. }

    % Step 3.3
    \solutionStep{6}{Apply the power rule}{ $\frac{d}{dx}(3x^2)=3\cdot2x^{2-1}=6x$ }

    % Step 3.4
    \solutionStep{6}{Differentiate the constant}{ $\frac{d}{dx}(7)=0$ }

    % Step 3.5
    \solutionStep{6}{Combine the results}{ $\frac{d}{dx}(3x^2+7)=6x$ }

    % Step 3.6
    \solutionStep{6}{Apply chain rule result}{ $\frac{d}{dx}\ln(3x^2+7)=\frac{6x}{3x^2+7}$ }

    % Step 4
    \mainStep{6}{Differentiate the denominator}

    % Step 4.1
    \solutionStep{6}{Apply chain rule again}{ $\frac{d}{dx}\ln(5x^2+6)=\frac{(5x^2+6)'}{5x^2+6}$ }

    % Step 4.2
    \solutionStep{6}{Differentiate $5x^2$}{ $5\cdot2x^{2-1}=10x$ }

    % Step 4.3
    \solutionStep{6}{Differentiate the constant}{ $\frac{d}{dx}(6)=0$ }

    % Step 4.4
    \solutionStep{6}{Combine the results}{ $\frac{d}{dx}(5x^2+6)=10x$ }

    % Step 4.5
    \solutionStep{6}{Apply chain rule result}{ $\frac{d}{dx}\ln(5x^2+6)=\frac{10x}{5x^2+6}$ }

    % Step 5
    \mainStep{6}{Rewrite the limit using derivatives}

    % Step 5.1
    \solutionStep{6}{Substitute derivatives into the limit}{ 
    $\displaystyle\lim_{x\to\infty}\frac{\frac{6x}{3x^2+7}}{\frac{10x}{5x^2+6}}$ }

    % Step 5.2
    \solutionStep{6}{Rewrite complex fraction as multiplication}{ 
    $\displaystyle\lim_{x\to\infty}\left(\frac{6x}{3x^2+7}\cdot\frac{5x^2+6}{10x}\right)$ }

    % Step 5.3
    \solutionStep{6}{Multiply numerators and denominators}{
    $\displaystyle\lim_{x\to\infty}\frac{6x(5x^2+6)}{10x(3x^2+7)}$ }

    % Step 5.4
    \solutionStep{6}{Distribute $6x$ in numerator}{
    $6x(5x^2+6)=30x^3+36x$ }

    % Step 5.5
    \solutionStep{6}{Distribute $10x$ in denominator}{
    $10x(3x^2+7)=30x^3+70x$ }

    % Step 5.6
    \solutionStep{6}{Rewrite the limit}{
    $\displaystyle\lim_{x\to\infty}\frac{30x^3+36x}{30x^3+70x}$ }

    % Step 6
    \mainStep{6}{Evaluate the limit as $x\to\infty$}

    % Step 6.1
    \solutionStep{6}{Factor highest power $x^3$ from numerator and denominator}{
    $\displaystyle\frac{x^3(30+\frac{36}{x^2})}{x^3(30+\frac{70}{x^2})}$ }

    % Step 6.2
    \solutionStep{6}{Cancel $x^3$}{
    $\displaystyle\frac{30+\frac{36}{x^2}}{30+\frac{70}{x^2}}$ }

    % Step 6.3
    \solutionStep{6}{Evaluate as $x\to\infty$}{
    $\frac{36}{x^2}\to0$ and $\frac{70}{x^2}\to0$ }

    % Step 6.4
    \solutionStep{6}{Compute final value}{
    $\displaystyle\frac{30+0}{30+0}=1$ }

    % Step 7
    \mainStep{6}{Final Answer: the correct option is B. 1 }

\end{tcolorbox}
\newpage

% QUESTION: 7
% STATUS: DRAFT
% AUTHOR: DG
% CREATED: 2026-02-02
% LAST UPDATED: 2026-02-09
% NOTES: 
% Q: If we perform the exact same operation twice, do we need to explain it in the same level of detail the second time? For example, if we showed how to apply L’Hôpital’s Rule in detail in Step 2, do we need to repeat the full explanation again in Step 4 when we apply L’Hôpital’s Rule a second time?
% Step 4 needs to add more steps.
% Maybe split Step 2 into 2 big steps. Right now it has 12 substeps, which is too long, and the substeps need to be further expand.

% Question
\section*{Question 7}
Evaluate $\displaystyle\lim_{x\to 0}\frac{1-\cos 2x}{x^3+5x^2}$.
    
% Multiple Choices
\begin{enumerate}[label = \alph*.]
    \item 0
    \item $\displaystyle\frac{1}{10}$
    \item $\displaystyle\frac{2}{5}$
    \item I don't know
\end{enumerate}

\setcounter{solutionMainStep}{0}
\setcounter{solutionSubStep}{0}

% Solution Box
\begin{tcolorbox}[breakable, colback= blue!5!white, colframe= blue!50!black, title=\textbf{Solution: c) $\displaystyle\frac{2}{5}$}]

    % Step 1
    \mainStep{7}{Identify the limit form}

    % Step 1.1
    \solutionStep{7}{Evaluate numerator at $x=0$}{
    \begin{flushleft}
    \begin{tabbing}
    \hspace{4.5cm} \= \kill
    $\begin{aligned}[t]
    1-\cos 2x &= 1-\cos(2\cdot0) \\
              &= 1-\cos(0) \\
              &= 1-1 \\
              &= 0
    \end{aligned}$
    \end{tabbing}
    \end{flushleft}
    }

    % Step 1.2
    \solutionStep{7}{Evaluate denominator at $x=0$}{
    \begin{flushleft}
    \begin{tabbing}
    \hspace{4.5cm} \= \kill
    $\begin{aligned}[t]
    x^3+5x^2 &= (0)^3+5(0)^2 \\
             &= 0+0 \\
             &= 0
    \end{aligned}$
    \end{tabbing}
    \end{flushleft}
    }

    % Step 1.3
    \solutionStep{7}{Determine indeterminate form}{
    $\lim_{x\to 0}\frac{1-\cos 2x}{x^3+5x^2} = \frac{0}{0}$}

    % Step 2
    \mainStep{7}{Apply L'Hôpital's Rule (First Time)}

    % Step 2.1
    \solutionStep{7}{State L'Hôpital's Rule}{
    \colorbox{yellow}{L'Hôpital's Rule:} If $\lim_{x\to a}\frac{f(x)}{g(x)}$ gives $\frac{0}{0}$ or $\frac{\infty}{\infty}$, then 
    $\displaystyle\lim_{x\to a}\frac{f(x)}{g(x)}=\lim_{x\to a}\frac{f'(x)}{g'(x)}$, provided the new limit exists.
    }

    % Step 2.2
    \solutionStep{7}{Differentiate numerator}{
    $\frac{d}{dx}(1-\cos 2x)
    =
    \frac{d}{dx}(1)
    +
    \frac{d}{dx}(-\cos 2x)$
    }

    % Step 2.3
    \solutionStep{7}{Differentiate $1$}{
    $\frac{d}{dx}(1)=0$
    }

    % Step 2.4
    \solutionStep{7}{Differentiate $-\cos 2x$ using chain rule}{
    \colorbox{yellow}{Chain Rule:} If $F(x)=\cos(h(x))$, then $\frac{d}{dx}[\cos(h(x))]=-\sin(h(x))h'(x)$.
    }

    % Step 2.5
    \solutionStep{7}{Differentiate inside function}{
    $\frac{d}{dx}(2x)=2$
    }

    % Step 2.6
    \solutionStep{7}{Combine chain rule result}{
    $\frac{d}{dx}(-\cos 2x)=(-1)(-\sin 2x)(2)=2\sin 2x$
    }

    % Step 2.7
    \solutionStep{7}{Combine numerator derivatives}{
    $\frac{d}{dx}(1-\cos 2x)=0+2\sin 2x=2\sin 2x$
    }

    % Step 2.8
    \solutionStep{7}{Differentiate denominator}{
    $\frac{d}{dx}(x^3+5x^2)
    =
    \frac{d}{dx}(x^3)
    +
    \frac{d}{dx}(5x^2)$
    }

    % Step 2.9
    \solutionStep{7}{Apply power rule}{
    \colorbox{yellow}{Power Rule:} If $g(x)=x^n$, then $g'(x)=nx^{n-1}$.
    }

    % Step 2.10
    \solutionStep{7}{Differentiate $x^3$}{
    $\frac{d}{dx}(x^3)=3x^2$
    }

    % Step 2.11
    \solutionStep{7}{Differentiate $5x^2$}{
    $\frac{d}{dx}(5x^2)=10x$
    }

    % Step 2.12
    \solutionStep{7}{Combine denominator derivatives}{
    $\frac{d}{dx}(x^3+5x^2)=3x^2+10x$
    }

    % Step 2.13
    \solutionStep{7}{Rewrite limit}{
    $\displaystyle\lim_{x\to0}\frac{2\sin 2x}{3x^2+10x}$
    }

    % Step 3
    \mainStep{7}{Check the new limit form}

    % Step 3.1
    \solutionStep{7}{Evaluate numerator at $x=0$}{
    \begin{flushleft}
    \begin{tabbing}
    \hspace{4.5cm} \= \kill
    $\begin{aligned}[t]
    2\sin 2x &= 2\sin(2\cdot0) \\
             &= 2\sin(0) \\
             &= 2(0) \\
             &= 0
    \end{aligned}$
    \end{tabbing}
    \end{flushleft}
    }

    % Step 3.2
    \solutionStep{7}{Evaluate denominator at $x=0$}{
    \begin{flushleft}
    \begin{tabbing}
    \hspace{4.5cm} \= \kill
    $\begin{aligned}[t]
    3x^2+10x &= 3(0)^2+10(0) \\
             &= 0+0 \\
             &= 0
    \end{aligned}$
    \end{tabbing}
    \end{flushleft}
    }

    % Step 3.3
    \solutionStep{7}{Determine indeterminate form}{
    $\lim_{x\to0}\frac{2\sin 2x}{3x^2+10x} = \frac{0}{0}$
    }

    % Step 4
    \mainStep{7}{Apply L'Hôpital's Rule (Second Time)}

    % Step 4.1
    \solutionStep{7}{Differentiate numerator}{
    $\frac{d}{dx}(2\sin 2x)=2(\cos 2x)(2)=4\cos 2x$
    }

    % Step 4.2
    \solutionStep{7}{Differentiate denominator}{
    $\frac{d}{dx}(3x^2+10x)=6x+10$
    }

    % Step 4.3
    \solutionStep{7}{Rewrite limit}{
    $\displaystyle\lim_{x\to0}\frac{4\cos 2x}{6x+10}$
    }

    % Step 5
    \mainStep{7}{Evaluate the limit}

    % Step 5.1
    \solutionStep{7}{Evaluate numerator at $x=0$}{
    $4\cos 2x=4\cos(2\cdot0)=4\cos(0)=4(1)=4$
    }

    % Step 5.2
    \solutionStep{7}{Evaluate denominator at $x=0$}{
    $6x+10=6(0)+10=10$
    }

    % Step 5.3
    \solutionStep{7}{Combine the result}{
    $\lim_{x\to0}\frac{4\cos 2x}{6x+10}=\frac{4}{10}=\frac{2}{5}$
    }

    % Step 6
    \mainStep{7}{Final Answer: the correct option is C. $\displaystyle\frac{2}{5}$ }

\end{tcolorbox}
\newpage

% QUESTION: 8
% STATUS: DRAFT
% AUTHOR: DG
% CREATED: 2026-02-02
% LAST UPDATED: 2026-02-09
% NOTES: 
% Need to expand Step 1.1 and 1.2
% For Step 3.4, Do we need to show how to repeat 25 times? 

% Question
\section*{Question 8}
Evaluate $\displaystyle\lim_{x\to\infty}\frac{x^{25}}{2^x}$.
    
% Multiple Choices
\begin{enumerate}[label = \alph*.]
    \item 0
    \item $\displaystyle\frac{25!}{(\ln 2)^{25}}$
    \item $\infty$
    \item I don't know
\end{enumerate}

\setcounter{solutionMainStep}{0}
\setcounter{solutionSubStep}{0}

% Solution Box
\begin{tcolorbox}[breakable, colback= blue!5!white, colframe= blue!50!black, title=\textbf{Solution: a) 0}]

    % Step 1
    \mainStep{8}{Identify the limit form}

    % Step 1.1
    \solutionStep{8}{Evaluate numerator as $x\to\infty$}{
    $x^{25}\to\infty$
    }

    % Step 1.2
    \solutionStep{8}{Evaluate denominator as $x\to\infty$}{
    $2^x\to\infty$
    }

    % Step 1.3
    \solutionStep{8}{Determine indeterminate form}{
    $\lim_{x\to\infty}\frac{x^{25}}{2^x}=\frac{\infty}{\infty}$}

    % Step 2
    \mainStep{8}{Apply L'Hôpital's Rule (First Time)}

    % Step 2.1
    \solutionStep{8}{L'Hôpital's Rule}{
    \colorbox{yellow}{L'Hôpital's Rule:} If $\lim_{x\to a}\frac{f(x)}{g(x)}$ gives $\frac{0}{0}$ or $\frac{\infty}{\infty}$, then 
    $\displaystyle\lim_{x\to a}\frac{f(x)}{g(x)}=\lim_{x\to a}\frac{f'(x)}{g'(x)}$.
    }

    % Step 2.2
    \solutionStep{8}{Differentiate numerator}{
    \colorbox{yellow}{Power Rule:} If $g(x)=x^n$, then $g'(x)=nx^{n-1}$.
    }

    % Step 2.3
    \solutionStep{8}{Apply power rule}{
    $\frac{d}{dx}(x^{25})=25x^{24}$
    }

    % Step 2.4
    \solutionStep{8}{Differentiate denominator}{
    \colorbox{yellow}{Exponential Derivative Rule:} $\frac{d}{dx}(a^x)=a^x\ln a$.
    }

    % Step 2.5
    \solutionStep{8}{Apply exponential rule}{
    $\frac{d}{dx}(2^x)=2^x\ln 2$
    }

    % Step 2.6
    \solutionStep{8}{Rewrite limit}{
    $\displaystyle\lim_{x\to\infty}\frac{25x^{24}}{2^x\ln 2}$
    }

    % Step 3
    \mainStep{8}{Apply L'Hôpital's Rule (Second Time)}

    % Step 3.1
    \solutionStep{8}{Differentiate numerator again}{
    $\frac{d}{dx}(25x^{24})=25\cdot24x^{23}$
    }

    % Step 3.2
    \solutionStep{8}{Differentiate denominator again}{
    $\frac{d}{dx}(2^x\ln 2)=\ln 2\cdot(2^x\ln 2)=2^x(\ln 2)^2$
    }

    % Step 3.3
    \solutionStep{8}{Rewrite limit}{
    $\displaystyle\lim_{x\to\infty}\frac{25\cdot24x^{23}}{2^x(\ln 2)^2}$
    }

    % Step 3.4
    \solutionStep{8}{Continue applying L'Hôpital}{
    After repeating this process 25 times
    }

    % Step 3.5
    \solutionStep{8}{Final differentiated form}{
    $\displaystyle\lim_{x\to\infty}\frac{25!}{2^x(\ln 2)^{25}}$
    }

    % Step 4
    \mainStep{8}{Evaluate the limit}

    % Step 4.1
    \solutionStep{8}{Analyze numerator}{
    $25!$ is a constant.
    }

    % Step 4.2
    \solutionStep{8}{Analyze denominator}{
    $2^x(\ln 2)^{25}\to\infty$
    }

    % Step 4.3
    \solutionStep{8}{Evaluate limit}{
    $\lim_{x\to\infty}\frac{25!}{2^x(\ln 2)^{25}} = \frac{\text{constant}}{\infty}=0$
    }

    % Step 5
    \mainStep{8}{Final Answer: the correct option is A. 0 }

\end{tcolorbox}
\newpage

\end{document}
